\documentclass[12pt]{article}
\title{\textbf{MATH 3670 Homework 1}}
\author{Saahas Swaroop}

\usepackage{latexsym}
\usepackage{hyperref}
\usepackage{amsmath}

\begin{document}
  \maketitle
  \pagebreak
  

  \section*{1}
  Sample space with replacement: $$ S = \{rr, rg, rb, gg, gr, gb, bb, br, bg\} $$
  Sample space without replacement: $$ S = \{rg, rb ,gr, gb, br, bg\} $$

  $r$: picking a red marble
  $g$: picking a green marble
  $b$: picking a blue marble
  $rg$: picking a red marble then a green marble

  \section*{2}
  Sample space: $$S = \{HHH, HHT, HTH, HTT, TTT, TTH, THT, THH\}$$
  Event: $$E = \{HHH, HHT, HTH, THH\}$$

  \section*{3}
  \subsection*{a}
  $\{7\}$
  \subsection*{b}
  $\{1, 3, 4, 5, 7\}$
  \subsection*{c}
  $\{3, 5, 7\}$
  \subsection*{d}
  $\{1, 3, 4, 5\}$
  \subsection*{e}
  $\{4, 6\}$
  \subsection*{f}
  $\{4, 7\}$

  \section*{4}
  $$ E = \{12, 14, 16, 21, 23, 25, 32, 34, 36, 41, 43, 45, 52, 54, 56, 61, 63, 65\} $$
  $$ F = \{11, 12, 13, 14, 15, 16\} $$
  $$ G = \{14, 23, 32, 41\} $$

  $$ EF = \{12, 14, 16\} $$
  $$ E \cup F = \{11, 12, 13, 14, 15, 16, 21, 23, 25, 32, 34, 36, 41, 43, 45, 52, 54, 56, 61, 63, 65\} $$
  $$ FG = \{14\} $$
  $$ EF^c = \{21, 23, 25, 32, 34, 36, 41, 43, 45, 52, 54, 56, 61, 63, 65\} $$
  $$ EG \cup FG = \{14\} $$

  \section*{7}
  \subsection*{a}
  $S$
  \subsection*{b}
  $\emptyset$
  \subsection*{c}
  $\emptyset$
  \subsection*{d}
  $\emptyset$
  \subsection*{e}
  $F \cup (E \cap G)$

  \section*{10}
  Let $G$ be an event such that $F = E \cup G$ and $EG = \emptyset$. This means that the cardinality of $F$ must be greater than or equal to that of $E$, as the smallest possible set $G$ is $\emptyset$. Assuming all events have equal probability, $P(E) \leq P(F)$ due to their respective cardinalities.

  \section*{14}
  The probability that exactly one of the events $E$ or $F$ occur is $P(E \cup F) - P(EF)$. Since $P(E \cup F) = P(E) + P(F) - P(EF)$, the original expression can be rewritten as $P(E) + P(F) - 2P(EF)$.

  \section*{16}
  $\binom{n}{r} = \frac{n!}{r!(n-r)!}$ and $\binom{n}{n-r}$ can be rewritten as $\frac{n!}{(n-r)!(n-(n-r))!}$ which can be simplified to $\frac{n!}{n!(n-r)!}$, which shows both forms are equal.

  An alternative method of proving this is by conceptually comparing what it means to choose $r$ items from a group of $n$ items. If we select $r$ items, there will be $n-r$ items left behind. Instead of selecting $r$ items in this way, we can instead select $n-r$ items, leaving $r$ items behind. These are both equivalent, showcasing how $\binom{n}{n-r}$ and $\binom{n}{r}$ are equivalent.

  \section*{18}
  \subsection*{a}
  $\frac{1}{3}$
  \subsection*{b}
  $\frac{1}{3}$
  \subsection*{c}
  $\frac{1}{15}$

  \section*{23}
  $$ P(R \ R) = \frac{P(RR)}{P(R)} = \frac{\frac{1}{3}}{\frac{2}{3}} = \frac{1}{2}$$


\end{document}

